% Figure colours remain semantic in both review and clean manuscripts.
\begingroup
\PaperShowRevisionsfalse
% Shared vector-figure vocabulary for the Chinese IEEE manuscript.
% Every figure in this directory inputs this file, so the manuscript only
% needs the standard TikZ package and the libraries already loaded by
% paper_zh.tex.  Keep the guard: figures may be input more than once while
% editing or when a stand-alone smoke-test document is built.
\ifdefined\GALKFigureStyleLoaded\else
\def\GALKFigureStyleLoaded{1}

% Baseline-derived visual tokens: ICCAD uses a restrained blue/orange device
% palette, while ZAC uses conventional black axes and a small fixed plot set.
\definecolor{galkInk}{HTML}{231F20}
\definecolor{galkMuted}{HTML}{666666}
\definecolor{galkGrid}{HTML}{D6D6D6}
\definecolor{galkPaper}{HTML}{FFFFFF}
% Semantic colors are fixed across every figure: orange=entanglement zone,
% blue=storage zone, black=atom identity, red=invalidity, and blue dashes=future
% prediction.  Plot and search-group colours are deliberately independent of
% the two physical-zone colours.
\definecolor{galkStorage}{HTML}{0064BC}
\definecolor{galkStorageFill}{HTML}{E8F1F7}
\definecolor{galkEntangle}{HTML}{E37323}
\definecolor{galkEntangleFill}{HTML}{F9ECE3}
\definecolor{galkResident}{HTML}{485CB3}
\definecolor{galkResidentFill}{HTML}{EEF0F8}
\definecolor{galkGood}{HTML}{A4AF07}
\definecolor{galkBad}{HTML}{C84E45}
\definecolor{galkLook}{HTML}{0064BC}
\definecolor{galkData}{HTML}{315E8A}
\definecolor{galkTail}{HTML}{B44B27}
\definecolor{galkGroupOne}{HTML}{4C78A8}
\definecolor{galkGroupTwo}{HTML}{7A7A7A}
\definecolor{galkGroupThree}{HTML}{9C755F}

\tikzset{
  % Shared engineering-figure contract.  Physical snapshots must place an
  % orange entanglement band above a blue storage field; only their scale may
  % change.  Labels are aligned at the upper-left of both regions.
  galk/base/.style={
    font=\rmfamily\fontsize{9.0pt}{10.0pt}\selectfont,
    text=galkInk,
    line cap=butt,
    line join=miter,
    >=Stealth
  },
  galk/panel title/.style={
    font=\rmfamily\bfseries\fontsize{9.2pt}{10.2pt}\selectfont,
    text=galkInk,
    anchor=west
  },
  galk/label/.style={
    font=\rmfamily\fontsize{9.0pt}{10.0pt}\selectfont,
    text=galkInk
  },
  galk/small label/.style={
    font=\rmfamily\fontsize{8.7pt}{9.7pt}\selectfont,
    text=galkInk
  },
  galk/zone label/.style={
    font=\rmfamily\bfseries\fontsize{8.7pt}{9.7pt}\selectfont,
    text=galkInk,
    anchor=north west,
    inner sep=.5pt
  },
  galk/entanglement zone/.style={
    draw=galkEntangle,
    fill=galkEntangleFill,
    line width=.40pt
  },
  galk/storage zone/.style={
    draw=galkStorage,
    fill=galkStorageFill,
    line width=.40pt
  },
  galk/site/.style={
    circle,
    draw=galkInk!68,
    fill=white,
    line width=.36pt,
    minimum size=1.35mm,
    inner sep=0pt
  },
  galk/occupied atom/.style={
    circle,
    draw=galkInk,
    fill=galkInk,
    line width=.36pt,
    minimum size=1.55mm,
    inner sep=0pt
  },
  galk/target site/.style={
    circle,
    draw=galkInk!60,
    fill=galkInk!10,
    line width=.34pt,
    minimum size=1.55mm,
    inner sep=0pt
  },
  galk/q label/.style={
    font=\rmfamily\itshape\fontsize{8.5pt}{9.5pt}\selectfont,
    text=galkInk,
    inner sep=.25pt
  },
  galk/rydberg pair/.style={
    draw=galkInk!62,
    densely dashed,
    line width=.36pt
  },
  galk/accepted path/.style={
    -{Stealth[length=.95mm,width=.62mm]},
    draw=galkInk,
    line width=.48pt
  },
  galk/alternative path/.style={
    -{Stealth[length=.95mm,width=.62mm]},
    draw=galkMuted,
    densely dashed,
    line width=.48pt
  },
  galk/future path/.style={
    -{Stealth[length=.95mm,width=.62mm]},
    draw=galkLook,
    densely dashed,
    line width=.48pt
  },
  galk/invalid path/.style={
    -{Stealth[length=.95mm,width=.62mm]},
    draw=galkBad,
    densely dashed,
    line width=.50pt
  }
}
\fi

\providecommand{\PaperRevision}[1]{#1}
\begin{tikzpicture}[
  galk/base,x=1cm,y=1cm,
  title/.style={font=\rmfamily\bfseries\fontsize{9pt}{10pt}\selectfont,align=center},
  label/.style={font=\rmfamily\fontsize{8.7pt}{9.7pt}\selectfont},
  auxiliary/.style={font=\rmfamily\fontsize{8.5pt}{9.5pt}\selectfont},
  qlabel/.style={galk/q label,font=\rmfamily\fontsize{8.5pt}{9.5pt}\selectfont},
  wire/.style={draw=galkInk!75,line width=.42pt},
  divider/.style={draw=galkInk!25,line width=.38pt},
  flow/.style={galk/accepted path,line width=.55pt},
  gate/.style={draw=galkInk!75,fill=white,line width=.42pt,
    minimum width=3.7mm,minimum height=2.8mm,inner sep=.3pt,
    font=\rmfamily\fontsize{8.5pt}{9.5pt}\selectfont},
  stage/.style={draw=galkInk!45,fill=white,line width=.45pt,
    minimum height=5.84cm,inner sep=0pt},
  module/.style={draw=galkInk!65,fill=white,line width=.45pt,
    align=center,inner sep=2pt,font=\rmfamily\fontsize{8.5pt}{9.5pt}\selectfont},
  trap/.style={galk/site},
  atom/.style={galk/occupied atom},
  feedback/.style={galk/future path}
]
  \path[use as bounding box] (0,0) rectangle (17.80,7.10);

  % One connected framework: each stage owns its visible inputs and outputs.
  % The complete current-layer decision/evaluation detail is inside joint.
  \node[stage,minimum width=2.89cm] (input) at (1.485,3.94) {};
  \node[stage,minimum width=3.02cm] (schedule) at (4.72,3.94) {};
  \node[stage,minimum width=2.02cm] (initial) at (7.52,3.94) {};
  \node[stage,minimum width=4.46cm,draw=galkInk,line width=.70pt] (joint) at (11.04,3.94) {};
  \node[stage,minimum width=2.06cm] (routing) at (14.58,3.94) {};
  \node[stage,minimum width=1.86cm] (output) at (16.82,3.94) {};
  \coordinate (flowRow) at (0,6.58);
  \foreach \from/\to in {input/schedule,schedule/initial,initial/joint,joint/routing,routing/output}
    \draw[flow] (\from.east |- flowRow)--(\to.west |- flowRow);
  \foreach \xx/\txt in {1.485/Input,4.72/Gate scheduling,11.04/Joint optimization,14.58/Routing,16.82/ZAIR}
    \node[title] at (\xx,6.58) {\PaperRevision{\txt}};
  \node[title] at (7.52,6.51)
    {\PaperRevision{Initial}\\\PaperRevision{placement}};

  % Input and scheduled circuits preserve four CZ gates and six 1Q gates.
  \def\czGate#1#2#3{%
    \draw[wire] (#1,#2)--(#1,#3);
    \fill[galkInk] (#1,#2) circle (.8pt);
    \node[gate] at (#1,#3) {\PaperRevision{Z}};
  }
  \foreach \qq/\yy in {0/5.84,1/5.34,2/4.84,3/4.34,4/3.84} {
    \node[qlabel,anchor=east] at (.31,\yy) {\PaperRevision{$q_{\qq}$}};
    \draw[wire] (.38,\yy)--(2.83,\yy);
    \node[qlabel,anchor=east] at (3.47,\yy) {\PaperRevision{$q_{\qq}$}};
    \draw[wire] (3.54,\yy)--(6.14,\yy);
  }
  \node[gate] at (.57,5.84) {\PaperRevision{RY}};
  \node[gate] at (.57,4.34) {\PaperRevision{RZ}};
  \czGate{.88}{5.84}{5.34}
  \czGate{1.23}{4.84}{4.34}
  \node[gate] at (1.59,5.34) {\PaperRevision{U}};
  \node[gate] at (1.59,3.84) {\PaperRevision{RY}};
  \czGate{1.91}{5.84}{4.84}
  \czGate{2.25}{5.34}{3.84}
  \node[gate] at (2.68,4.84) {\PaperRevision{RZ}};
  \node[gate] at (2.68,4.34) {\PaperRevision{U}};

  % q4 has no preceding 2Q gate; q3's U follows the first 2Q layer.
  \foreach \xx in {3.95,4.47,4.94,5.75}
    \draw[divider,densely dashed] (\xx,3.62)--(\xx,6.05);
  \node[gate] at (3.73,5.84) {\PaperRevision{RY}};
  \node[gate] at (3.73,4.34) {\PaperRevision{RZ}};
  \node[gate] at (3.73,3.84) {\PaperRevision{RY}};
  \czGate{4.22}{5.84}{5.34}
  \czGate{4.22}{4.84}{4.34}
  \node[gate] at (4.71,5.34) {\PaperRevision{U}};
  \node[gate] at (4.71,4.34) {\PaperRevision{U}};
  \czGate{5.20}{5.84}{4.84}
  \czGate{5.53}{5.34}{3.84}
  \node[gate] at (5.98,4.84) {\PaperRevision{RZ}};
  \foreach \xx/\lab in {3.73/{$1Q_0$},4.71/{$1Q_1$},5.98/{$1Q_2$}}
    \node[auxiliary] at (\xx,3.44) {\PaperRevision{\lab}};
  \foreach \xx/\lab in {4.22/{$2Q_0$},5.365/{$2Q_1$}}
    \node[auxiliary] at (\xx,3.12) {\PaperRevision{\lab}};
  \node[module,text width=2.44cm,minimum height=.78cm] (layerSequence)
    at (4.72,2.21)
    {\PaperRevision{Ordered 2Q layers\\+ associated 1Q gates}};
  \draw[flow] (4.72,2.92)--(layerSequence.north);

  % Architecture is an input, not an extra atom-placement example.
  \node[label] at (1.485,3.30) {\PaperRevision{Architecture}};
  \path[galk/entanglement zone] (.39,2.55) rectangle (2.62,3.02);
  \path[galk/storage zone] (.39,1.42) rectangle (2.62,2.40);
  \foreach \xx in {.64,.96,2.04,2.36}
    \node[trap] at (\xx,2.78) {};
  \foreach \xx in {.64,1.50,2.36}
    \foreach \yy in {1.61,1.89,2.17}
      \node[trap] at (\xx,\yy) {};
  \draw[galk/rydberg pair] (.80,2.78) ellipse (.25 and .15);
  \draw[galk/rydberg pair] (2.20,2.78) ellipse (.25 and .15);
  \node[auxiliary] at (.20,2.78) {\PaperRevision{$\mathcal E$}};
  \node[auxiliary] at (.20,1.91) {\PaperRevision{$\mathcal S$}};

  % Initialization is evaluated once, before the repeated layer loop.
  \node[module,text width=1.52cm,minimum height=.64cm] (initialCandidates)
    at (7.52,5.60) {\PaperRevision{Candidate\\layouts $\pi$}};
  \node[module,text width=1.52cm,minimum height=.82cm] (initialEval)
    at (7.52,4.40) {\PaperRevision{Evaluate\\early layers}};
  \node[module,text width=1.52cm,minimum height=.62cm] (initialMap)
    at (7.52,3.18) {\PaperRevision{Select layout}};
  \node[label,align=center] (initialState) at (7.52,1.98)
    {\PaperRevision{Starting state}\\\PaperRevision{$\mathbf{s}_0$}};
  \draw[flow] (initialCandidates.south)--(initialEval.north);
  \draw[flow] (initialEval.south)--(initialMap.north);
  \draw[flow] (initialMap.south)--(initialState.north);

  % Candidate evaluation belongs to the joint stage. The physical snapshot
  % is the fixed 2Q1 example, whereas the surrounding loop is generic.
  \node[module,text width=3.92cm,minimum height=.61cm] (jointSearch)
    at (11.04,5.94)
    {\PaperRevision{Joint search + matching\\Gate sites / residence / storage}};
  \node[draw=galkInk!40,line width=.38pt,minimum width=4.10cm,
    minimum height=2.47cm,inner sep=0pt] (currentEval) at (11.04,4.175) {};
  \node[auxiliary] at (11.04,5.18)
    {\PaperRevision{Example: evaluate $2Q_1+1Q_2$}};
  \draw[flow] (jointSearch.south)--(currentEval.north);

  \def\device#1#2{%
    \begin{scope}[shift={(#1,#2)},x=.80cm,y=.80cm]
      \path[galk/entanglement zone] (0,.98) rectangle (2.00,2.10);
      \path[galk/storage zone] (0,0) rectangle (2.00,.88);
      \foreach \xx in {.24,.56,1.44,1.76}
        \node[trap] at (\xx,1.55) {};
      \foreach \xx in {.24,1.00,1.76}
        \foreach \yy in {.20,.48,.76}
          \node[trap] at (\xx,\yy) {};
      \draw[galk/rydberg pair] (.40,1.55) ellipse (.26 and .17);
      \draw[galk/rydberg pair] (1.60,1.55) ellipse (.26 and .17);
      \node[auxiliary] at (.40,1.88) {\PaperRevision{$T_1$}};
      \node[auxiliary] at (1.60,1.88) {\PaperRevision{$T_2$}};
    \end{scope}%
  }
  \def\routeInitial{0/.24/1.55,1/.56/1.55,2/1.44/1.55,3/1.76/1.55,4/1.00/.20}
  \def\routeFinal{0/.24/1.55,1/1.44/1.55,2/.56/1.55,3/1.76/.48,4/1.76/1.55}
  \device{9.08}{3.04}
  \device{11.42}{3.04}
  \foreach \qq/\xx/\yy in \routeInitial
    \node[atom] at (9.08+.80*\xx,3.04+.80*\yy) {};
  \foreach \qq/\xx/\yy in \routeFinal
    \node[atom] at (11.42+.80*\xx,3.04+.80*\yy) {};
  \foreach \xx/\lab in {9.40/{$q_0,q_1$},10.36/{$q_2,q_3$},11.74/{$q_0,q_2$},12.70/{$q_1,q_4$}}
    \node[qlabel] at (\xx,4.01) {\PaperRevision{\lab}};
  \node[qlabel,anchor=west] at (10.02,3.20) {\PaperRevision{$q_4$}};
  \node[qlabel] at (12.55,3.20) {\PaperRevision{$q_3$}};
  \node[label] (stateBefore) at (9.88,4.87) {\PaperRevision{$\mathbf{s}_\ell$}};
  \node[label] at (12.22,4.87) {\PaperRevision{$\mathbf{s}_{\ell+1}$}};
  \coordinate (stateAfter) at (12.22,3.04);
  \draw[flow] (10.76,3.91)--(11.34,3.91);

  \node[module,draw=galkLook!60,text width=3.92cm,minimum height=1.40cm]
    (future) at (11.04,1.82) {};
  \node[auxiliary] at (11.04,2.34) {\PaperRevision{Multi-layer look-ahead}};
  \node[auxiliary] at (11.04,1.99) {\PaperRevision{If layers remain}};
  \node[auxiliary] (futureOne) at (9.66,1.66) {\PaperRevision{$L_{\ell+1}$}};
  \node[auxiliary] (futureTwo) at (10.99,1.66) {\PaperRevision{$L_{\ell+2}$}};
  \node[auxiliary] (futureMore) at (12.35,1.66) {\PaperRevision{$\cdots$}};
  \draw[feedback] (stateAfter)--(12.22,2.53);
  \draw[feedback] (futureOne.east)--(futureTwo.west);
  \draw[feedback] (futureTwo.east)--(futureMore.west);
  \node[auxiliary] at (11.04,1.31) {\PaperRevision{$J_\ell$: current + future loss}};
  \draw[feedback] (future.west)--(8.90,1.82)--(8.90,5.94)--(jointSearch.west);

  % Routing applies only the selected current-layer plan. Its resulting
  % positions, occupancy and clock drive the next actual optimization.
  \node[module,text width=1.60cm,minimum height=.69cm] (selectedPlan)
    at (14.58,5.71) {\PaperRevision{Selected plan}};
  \node[module,text width=1.60cm,minimum height=.72cm] (routeBatches)
    at (14.58,4.63) {\PaperRevision{AOD-safe\\batches}};
  \node[module,text width=1.60cm,minimum height=.83cm] (executeCurrent)
    at (14.58,3.44) {\PaperRevision{Execute current\\2Q + 1Q}};
  \node[module,text width=1.60cm,minimum height=.98cm] (stateUpdate)
    at (14.58,2.02) {\PaperRevision{Update positions,\\occupancy\\and clock}};
  \draw[flow] (selectedPlan.south)--(routeBatches.north);
  \draw[flow] (routeBatches.south)--(executeCurrent.north);
  \draw[flow] (executeCurrent.south)--(stateUpdate.north);
  \draw[flow] (stateUpdate.south)--(14.58,.52)--(8.67,.52)--(8.67,6.58)--(joint.west |- flowRow);
  \node[label,fill=white,inner sep=2pt] at (11.72,.52)
    {\PaperRevision{Next layer: updated state}};

  % Full ZAIR order for the shown two-2Q-layer circuit, including final 1Q2.
  \foreach \yy/\txt in {5.88/init,5.39/1qGate,4.90/rearrangeJob,
    4.41/rydberg,3.92/1qGate,3.43/rearrangeJob,2.94/rydberg,2.45/1qGate}
    \node[auxiliary] at (16.82,\yy) {\PaperRevision{\txt}};
  \node[label,align=center] at (16.82,1.56)
    {\PaperRevision{Executable}\\\PaperRevision{instructions}};
\end{tikzpicture}
\endgroup
